\subsection{Complete recognition gains and unsuccessful-search costs}

The benchmark compares the published raw-then-span stage with the new
four-trial default, except for the explicitly labeled 24-trial case. Each
timed endpoint parses the input PD, runs the complete recognizer and exact
fallback when needed, independently replays any returned native proof,
and hashes and counts its serialized certificate. Earlier shortcuts are
disabled to make the native stage reachable. Both arms use no local seed
cap and retain a twelve-second global deadline. The native stage's
ordinary disabled default is not being changed by this benchmark.

Runs are serial, without overlapping tests, audits, builds or commits.
Every round randomly orders old, old A/A, new and new A/A calls. The initial
five-round experiment completes all 160 measured and 32 warm-up calls in
16.851 seconds, but its unchanged raw-success control has a new-arm A/A
ratio of 1.237 and an old/new ratio of 0.863. That variation warrants a
repeat before interpreting the control. The full initial record remains
in \path{data/cocycle-trees-initial-recognize.json}.

The final nine-round experiment completes all 288 measured and 32 warm-up
calls in 28.078 seconds. The table gives median primary-arm times and
median paired ratios; these ratios need not equal ratios of median times.
All verdicts and each implementation's proof hashes agree across repeated
calls and both experiments. Existing positive proofs remain byte-identical
across implementations. Newly successful native calls intentionally replace
the old inconclusive stage followed by an exact fallback.

\begin{center}
\small\begin{tabular}{lrrrrr}
\toprule
Input & Old (ms) & New (ms) & Old/new & Old A/A & New A/A \\
\midrule
Raw circle & 65.061 & 63.657 & 1.026 & 1.025 & 1.005 \\
Optimized positive & 129.131 & 144.014 & 0.925 & 0.991 & 0.978 \\
Early trial 2 & 50.894 & 38.314 & 1.345 & 1.017 & 0.988 \\
Early trial 3 & 78.248 & 51.757 & 1.516 & 0.995 & 0.991 \\
Five-crossing early & 110.901 & 65.990 & 1.639 & 0.958 & 0.989 \\
Late, 24-trial option & 111.954 & 181.031 & 0.624 & 1.024 & 1.014 \\
Genus-one miss & 51.600 & 59.570 & 0.866 & 1.004 & 1.000 \\
Trefoil & 50.198 & 57.041 & 0.877 & 0.982 & 1.024 \\
\bottomrule
\end{tabular}

\end{center}

The three early successes improve complete recognition by
$1.35$--$1.64\times$. They are \code{random-99}, \code{random-29} and
\code{random-124} from the preceding table. The unchanged raw-success
control is near parity: its 1.026 paired ratio is comparable to its 1.025
old-arm A/A ratio. There is no measured algorithmic gain to claim for
that branch.

The unsuccessful prelude has a visible cost. The previous optimized
positive is about eight percent slower by the paired ratio, while the
genus-one miss and trefoil cost about 15 and 14 percent more respectively.
Both misses still use the exact Khovanov fallback and return their expected
\code{UNKNOT} and \code{KNOTTED} verdicts. The optional 24-trial example
produces a new native disc proof, but its full call is about 60 percent
slower than the old stage plus fallback. The initial run also shows this
cost, with a paired ratio of 0.540. This is additional native proof
coverage, not a speedup for that example; the longer search is therefore
an explicit option rather than the default allowance.

These measurements support a bounded cheap prelude with a documented
tradeoff. They do not establish a uniform speedup, a solved-corpus gain for
the ordinary portfolio, or a global complexity improvement. The zero-trial
option retains the earlier candidate schedule when its cost is preferable.
